% SHARD-ID: AQM-23-SPEC-Z6-RESIDUE-QUDIT-FACTORISATION
% SHARD-TITLE: Spec Z/6Z and Residue Qudit Factorisation
% SHARD-SUMMARY: Computes the prime ideals and residue fields of \(\operatorname{Spec}(\mathbb Z/6\mathbb Z)\).
% SHARD-SUMMARY: Builds the residue-field Weyl label group and proves it quantizes to qubit tensor qutrit.
% SHARD-SUMMARY: Separates the label direct sum from the Hilbert-space tensor product and rules out a Hilbert direct sum.
% SHARD-KEYWORDS: arithmetic quantum field, Spec, residue field, Z/6Z, qubit, qutrit, tensor product, Weyl-Heisenberg

\section{Spec Z/6Z and Residue Qudit Factorisation}
\label{sec:spec-z6-residue-qudit-factorisation}

\subsection{Status and Source Anchors}

This shard records the morally preferred convention suggested by
Section~\ref{sec:prime-field-arithmetic-fields-qudit-paulis}: for a field
\(\mathbb F_p\), the arithmetic point space is \(\operatorname{Spec}\mathbb F_p\),
so the local quantum system is one \(p\)-level qudit.  For a non-field finite
ring there is no single quotient field.  The concrete replacement used here is
one local phase space over each residue field.

The ring calculation for \(\mathbb Z/6\mathbb Z\) is proved directly below.
The Weyl/Pauli facts use the same local sources as AQM-22: Ashikhmin--Knill
\path{references/symplectic_qecc/ashikhmin_knill_0005008_source/n9.tex},
lines 249--313 for finite-field shift/phase bases and product laws, and
Bombin--Martin--Delgado
\path{references/toric_code/bombin_martin_delgado_0605094_source/HomologicalErrorCorrection6.tex},
lines 1208--1259 for qudit Paulis, symplectic commutators, and projective
label maps.

\subsection{Lamport Dependency Skeleton}

\[
\begin{array}{rcl}
D1&:&R=\mathbb Z/6\mathbb Z.\\
L1&:&R\cong\mathbb F_2\times\mathbb F_3\text{ by Chinese remainders.}\\
L2&:&\operatorname{Spec}R=\{\mathfrak p_2=(2),\mathfrak p_3=(3)\}.\\
D2&:&\kappa(\mathfrak p_2)=\mathbb F_2,\quad
       \kappa(\mathfrak p_3)=\mathbb F_3.\\
D3&:&E_R=\kappa(\mathfrak p_2)^2\oplus\kappa(\mathfrak p_3)^2.\\
D4&:&\Omega_R=\Omega_2\oplus\Omega_3,\text{ with each }\Omega_i
       \text{ the standard }zx'-xz'.\\
L3&:&E_R\text{ is an orthogonal direct sum of local symplectic groups.}\\
D5&:&\mathcal H_R=\ell^2(\mathbb F_2)\otimes\ell^2(\mathbb F_3).\\
L4&:&W_R(e_2,e_3)=W_2(e_2)\otimes W_3(e_3)
       \text{ has the product Weyl law.}\\
T1&:&\operatorname{span}\{W_R(e):e\in E_R\}
       \cong M_2(\mathbb C)\otimes M_3(\mathbb C)\cong M_6(\mathbb C).\\
T2&:&\text{The Hilbert space is a tensor product, not }
       \mathbb C^2\oplus\mathbb C^3.
\end{array}
\]

\subsection{The Spectrum L1--D2}

Let \(R=\mathbb Z/6\mathbb Z\).  Define
\[
  \Phi:R\longrightarrow\mathbb F_2\times\mathbb F_3,\qquad
  \Phi(\bar n)=(\bar n\bmod 2,\bar n\bmod 3).
\]

\paragraph{Lemma \(L1\).}
\(\Phi\) is a ring isomorphism.

\emph{Proof.}
It is a well-defined ring homomorphism because congruence modulo \(6\) implies
congruence modulo \(2\) and modulo \(3\).  Its kernel consists of residue
classes divisible by both \(2\) and \(3\), hence by \(6\), so the kernel is
zero.  Both rings have six elements, so an injective map is bijective.

\paragraph{Lemma \(L2\).}
\[
  \operatorname{Spec}R=\{\mathfrak p_2,\mathfrak p_3\},
  \qquad
  \mathfrak p_2=(2)=\{0,2,4\},\quad
  \mathfrak p_3=(3)=\{0,3\}.
\]
Moreover \(R/\mathfrak p_2\cong\mathbb F_2\) and
\(R/\mathfrak p_3\cong\mathbb F_3\).

\emph{Proof.}
Under \(\Phi\), the ideal \(\mathfrak p_2=(2)\) maps to
\(0\times\mathbb F_3\), and \(\mathfrak p_3=(3)\) maps to
\(\mathbb F_2\times0\).  The quotients are therefore \(\mathbb F_2\) and
\(\mathbb F_3\), respectively, so both ideals are prime.

Every ideal \(J\) of \(\mathbb F_2\times\mathbb F_3\) is \(I_2\times I_3\):
if \((a,b)\in J\), then \((a,0)=(a,b)(1,0)\in J\) and
\((0,b)=(a,b)(0,1)\in J\).  Thus the proper ideals are \(0\times0\),
\(0\times\mathbb F_3\), and \(\mathbb F_2\times0\).  The quotient by
\(0\times0\) is \(\mathbb F_2\times\mathbb F_3\), which is not a domain
because \((1,0)(0,1)=0\).  Hence \(0\times0\) is not prime, and the two ideals
listed above are the only prime ideals.

The residue fields are
\[
  \kappa(\mathfrak p_2)=R/\mathfrak p_2\cong\mathbb F_2,\qquad
  \kappa(\mathfrak p_3)=R/\mathfrak p_3\cong\mathbb F_3.
\]

\subsection{Residue-Field Phase Space D3--L3}

For a point \(\mathfrak p\), define its local phase space by
\[
  V_{\mathfrak p}=\kappa(\mathfrak p)_X\oplus\kappa(\mathfrak p)_Z,
  \qquad
  \Omega_{\mathfrak p}((x,z),(x',z'))=zx'-xz'.
\]
For \(R=\mathbb Z/6\mathbb Z\), the global label group is
\[
  E_R=V_{\mathfrak p_2}\oplus V_{\mathfrak p_3}
     =\mathbb F_2^2\oplus\mathbb F_3^2 .
\]
This is a direct sum of finite abelian groups, not a vector space over one
field.  Define the alternating bicharacter
\[
  B_R(e,f)
  =
  (-1)^{\Omega_{\mathfrak p_2}(e_2,f_2)}
  \exp\!\left(\frac{2\pi i}{3}\Omega_{\mathfrak p_3}(e_3,f_3)\right),
\]
where \(e=(e_2,e_3)\) and \(f=(f_2,f_3)\).

\paragraph{Lemma \(L3\).}
The two local summands are orthogonal and nondegenerate: if
\(B_R(e,f)=1\) for all \(f\in E_R\), then \(e=0\).

\emph{Proof.}
Orthogonality is built into the formula: \(e_2\) only pairs with \(f_2\), and
\(e_3\) only pairs with \(f_3\).  The local form is nondegenerate because a
nonzero \((x,z)\) pairs nontrivially with either \((1,0)\) or \((0,1)\).  If
\(B_R(e,f)=1\) for all \(f\), first take \(f_3=0\), forcing \(e_2=0\).  Then
take \(f_2=0\), forcing \(e_3=0\).  Thus \(e=0\).

\subsection{Weyl Representation L4--T1}

Let
\[
  \mathcal H_2=\ell^2(\mathbb F_2),\qquad
  \mathcal H_3=\ell^2(\mathbb F_3),\qquad
  \mathcal H_R=\mathcal H_2\otimes\mathcal H_3.
\]
For \(i=2,3\), let \(W_i(x,z)=T_i(x)R_i(z)\) be the usual local
shift-phase Weyl operator on \(\ell^2(\mathbb F_i)\).  Define
\[
  W_R(e_2,e_3)=W_2(e_2)\otimes W_3(e_3).
\]

\paragraph{Lemma \(L4\).}
For \(e=(e_2,e_3)\) and \(f=(f_2,f_3)\),
\[
  W_R(e)W_R(f)=c_R(e,f)W_R(e+f),
\]
where
\[
  c_R(e,f)=(-1)^{z_2x'_2}
  \exp\!\left(\frac{2\pi i}{3}z_3x'_3\right)
\]
if \(e_i=(x_i,z_i)\) and \(f_i=(x'_i,z'_i)\).  Consequently
\[
  W_R(e)W_R(f)=B_R(e,f)W_R(f)W_R(e).
\]

\emph{Proof.}
The local Weyl laws give
\[
  W_i(e_i)W_i(f_i)=c_i(e_i,f_i)W_i(e_i+f_i).
\]
Tensoring the two equations gives the displayed product law with
\(c_R=c_2c_3\).  Dividing the product law for \(e,f\) by the one for \(f,e\)
gives the displayed commutator bicharacter.

\paragraph{Theorem \(T1\).}
\[
  \operatorname{span}_{\mathbb C}\{W_R(e):e\in E_R\}
  =
  \operatorname{End}(\mathcal H_2)\otimes\operatorname{End}(\mathcal H_3)
  \cong M_2(\mathbb C)\otimes M_3(\mathbb C)
  \cong M_6(\mathbb C).
\]

\emph{Proof.}
For a single \(p\)-level qudit, the local source records that the \(p^2\)
operators \(T^xR^z\) form an operator basis.  Thus
\(\{W_2(e_2)\}\) is a basis of \(M_2(\mathbb C)\), and \(\{W_3(e_3)\}\) is a
basis of \(M_3(\mathbb C)\).  Tensor products of bases form a basis of the
tensor product vector space, so the \(4\cdot9=36\) operators \(W_R(e_2,e_3)\)
span \(M_2(\mathbb C)\otimes M_3(\mathbb C)\).  The last isomorphism is the
standard Kronecker product identification.

\subsection{Direct Sum or Tensor Product? T2}

\paragraph{Theorem \(T2\).}
The quantum system attached to \(\operatorname{Spec}(\mathbb Z/6\mathbb Z)\)
by the residue-field Weyl rule is a qubit tensor a qutrit,
\[
  \mathcal H_R\cong\mathbb C^2\otimes\mathbb C^3\cong\mathbb C^6,
\]
not the Hilbert direct sum \(\mathbb C^2\oplus\mathbb C^3\).

\emph{Proof.}
The direct sum occurs in the label group:
\[
  E_R=\mathbb F_2^2\oplus\mathbb F_3^2.
\]
Weyl quantization sends an orthogonal direct sum of independent phase spaces
to tensor products of the local Weyl representations, as shown explicitly in
Lemma \(L4\).

A Hilbert direct sum would have operator algebra
\[
  \operatorname{End}(\mathbb C^2)\oplus\operatorname{End}(\mathbb C^3)
  \cong M_2(\mathbb C)\oplus M_3(\mathbb C),
\]
whose complex dimension is \(4+9=13\) and whose center is two-dimensional.
The Weyl operators from \(E_R\) instead give \(36\) independent operators and
the full matrix algebra \(M_6(\mathbb C)\), whose center is one-dimensional.
Thus the direct-sum Hilbert interpretation is incompatible with the Weyl label
count and with Theorem \(T1\).

\subsection{Conclusion}

The expectation "qutrit plus qubit" is correct only at the level of local
factors.  The precise statement is
\[
  \operatorname{Spec}(\mathbb Z/6\mathbb Z)
  =
  \{\mathfrak p_2,\mathfrak p_3\},
  \qquad
  \kappa(\mathfrak p_2)=\mathbb F_2,\quad
  \kappa(\mathfrak p_3)=\mathbb F_3,
\]
so the local systems are one qubit and one qutrit.  The global Weyl label
space is their orthogonal direct sum, but the Hilbert space is their tensor
product.  Equivalently, the observable algebra is \(M_2\otimes M_3\cong M_6\),
not \(M_2\oplus M_3\).
